\documentclass[aspectratio=169,10pt]{beamer}

\usepackage[UTF8,fontset=fandol]{ctex}
\usepackage{amsmath,amssymb}
\usepackage{booktabs}
\usepackage{tabularx}
\usepackage{listings}
\usepackage{dlutbeamer}

\title[本科毕业设计答辩]{面向复杂场景的智能分析系统设计与实现}
\subtitle{大连理工大学本科生毕业设计（论文）答辩}
\author[张三]{张三}
\institute[大连理工大学]{大连理工大学\quad 软件学院}
\date[2026 年 6 月]{2026 年 6 月}
\SetDLUTTitleExtra{学号：2022123456\quad 指导教师：李老师\ 教授}

\lstset{
  basicstyle=\ttfamily\scriptsize,
  keywordstyle=\color{DLUTNavy}\bfseries,
  commentstyle=\color{DLUTMuted},
  stringstyle=\color{DLUTOrange},
  showstringspaces=false,
  frame=single,
  rulecolor=\color{DLUTLine},
  backgroundcolor=\color{DLUTWarm},
  breaklines=true
}

\begin{document}

\begin{frame}[plain]
  \titlepage
\end{frame}

\begin{frame}{目录}
  \tableofcontents
\end{frame}

\section{课题背景}

\begin{frame}{研究背景}
  \begin{columns}[T,onlytextwidth]
    \begin{column}{.61\textwidth}
      \begin{itemize}
        \item 复杂业务场景中的数据规模不断增长，传统人工分析流程难以兼顾效率与稳定性。
        \item 现有系统在\alert{实时性、可解释性与部署成本}之间仍存在明显权衡。
        \item 本课题面向可落地的本科毕业设计，完成需求分析、方法设计、系统实现与实验验证。
      \end{itemize}
      \vspace{.2cm}
      \begin{alertblock}{核心问题}
        如何在有限资源下实现准确、稳定、可维护的智能分析流程？
      \end{alertblock}
    \end{column}
    \begin{column}{.36\textwidth}
      \DLUTPlaceholder[4.65cm]{替换为场景图片、\\系统截图或问题示意图}
    \end{column}
  \end{columns}
\end{frame}

\begin{frame}{研究目标与主要工作}
  \begin{columns}[T,onlytextwidth]
    \begin{column}{.31\textwidth}
      \begin{block}{01\quad 需求建模}
        梳理关键业务流程，明确性能指标与约束条件。
      \end{block}
    \end{column}
    \begin{column}{.31\textwidth}
      \begin{block}{02\quad 方法设计}
        构建可解释的分析方法，并给出模块化实现方案。
      \end{block}
    \end{column}
    \begin{column}{.31\textwidth}
      \begin{block}{03\quad 实验验证}
        通过对比实验与消融实验验证系统有效性。
      \end{block}
    \end{column}
  \end{columns}
  \vspace{.35cm}
  \DLUTTag{关键词}\quad
  \DLUTTag{模块化设计}\quad
  \DLUTTag{性能优化}\quad
  \DLUTTag{实验评估}
\end{frame}

\section{方案设计}

\begin{frame}{总体技术路线}
  \centering
  \begin{tikzpicture}[
      node distance=1.0cm,
      box/.style={
        draw=DLUTLine,
        fill=DLUTWarm,
        rounded corners=3pt,
        minimum width=2.42cm,
        minimum height=1.06cm,
        align=center,
        text=DLUTDeep,
        font=\small\bfseries
      },
      arrow/.style={->,>=stealth,thick,draw=DLUTBlue}
    ]
    \node[box] (data) {数据采集\\与预处理};
    \node[box,right=of data] (method) {核心方法\\与特征分析};
    \node[box,right=of method] (service) {服务封装\\与接口设计};
    \node[box,right=of service] (eval) {实验评估\\与迭代优化};
    \draw[arrow] (data) -- (method);
    \draw[arrow] (method) -- (service);
    \draw[arrow] (service) -- (eval);
  \end{tikzpicture}
  \vspace{.62cm}
  \begin{exampleblock}{设计原则}
    以清晰的数据流组织模块边界，保证算法原型可以平滑迁移到系统实现。
  \end{exampleblock}
\end{frame}

\begin{frame}{核心方法}
  \begin{columns}[T,onlytextwidth]
    \begin{column}{.54\textwidth}
      \begin{itemize}
        \item 定义输入样本 $x$ 的表示向量 $\boldsymbol{h}(x)$；
        \item 通过可学习参数 $\theta$ 估计任务输出；
        \item 使用验证集选择关键超参数，并分析误差来源。
      \end{itemize}
      \vspace{.22cm}
      \begin{block}{优化目标}
        \[
          \theta^\star =
          \arg\min_{\theta}
          \frac{1}{N}\sum_{i=1}^{N}
          \mathcal{L}\bigl(f_{\theta}(x_i), y_i\bigr)
          + \lambda \lVert\theta\rVert_2^2
        \]
      \end{block}
    \end{column}
    \begin{column}{.42\textwidth}
      \DLUTPlaceholder[4.78cm]{替换为模型结构图、\\算法流程图或伪代码截图}
    \end{column}
  \end{columns}
\end{frame}

\begin{frame}{系统架构}
  \centering
  \begin{tikzpicture}[
      layer/.style={
        draw=DLUTLine,
        rounded corners=3pt,
        minimum width=10.6cm,
        minimum height=.78cm,
        align=center,
        text=DLUTDeep,
        font=\small
      }
    ]
    \node[layer,fill=DLUTSky!55] (ui) {交互层：可视化页面 \quad 结果检索 \quad 任务管理};
    \node[layer,fill=DLUTWarm,below=.27cm of ui] (api) {服务层：REST API \quad 权限控制 \quad 日志记录};
    \node[layer,fill=DLUTSky!55,below=.27cm of api] (core) {算法层：数据预处理 \quad 核心模型 \quad 指标计算};
    \node[layer,fill=DLUTWarm,below=.27cm of core] (data) {数据层：业务数据 \quad 配置文件 \quad 实验记录};
  \end{tikzpicture}
\end{frame}

\begin{frame}[fragile]{关键模块实现}
  \begin{columns}[T,onlytextwidth]
    \begin{column}{.54\textwidth}
\begin{lstlisting}[language=Python]
def analyze(sample, model, threshold=0.5):
    features = preprocess(sample)
    score = model.predict(features)
    return {
        "score": float(score),
        "accepted": score >= threshold,
    }
\end{lstlisting}
    \end{column}
    \begin{column}{.42\textwidth}
      \begin{itemize}
        \item 统一输入输出格式，降低模块耦合度；
        \item 对关键参数提供可配置入口；
        \item 保存中间结果，方便复现实验与定位问题。
      \end{itemize}
    \end{column}
  \end{columns}
\end{frame}

\section{实验与结果}

\begin{frame}{实验设置}
  \begin{columns}[T,onlytextwidth]
    \begin{column}{.54\textwidth}
      \begin{table}
        \centering
        \small
        \begin{tabularx}{\linewidth}{>{\bfseries}lX}
          \toprule
          项目 & 设置 \\
          \midrule
          数据划分 & 训练集 / 验证集 / 测试集 = 7 : 1 : 2 \\
          评价指标 & Accuracy、F1-score、平均响应时间 \\
          对比方法 & Baseline、Method A、本文方法 \\
          运行环境 & Python 3.11，单卡 GPU \\
          \bottomrule
        \end{tabularx}
      \end{table}
    \end{column}
    \begin{column}{.42\textwidth}
      \begin{block}{实验关注点}
        \begin{itemize}
          \item 整体性能是否提升；
          \item 关键模块是否有效；
          \item 系统开销是否可接受。
        \end{itemize}
      \end{block}
    \end{column}
  \end{columns}
\end{frame}

\begin{frame}{实验结果}
  \begin{columns}[T,onlytextwidth]
    \begin{column}{.54\textwidth}
      \centering
      \begin{tikzpicture}[x=.95cm,y=.048cm]
        \draw[->,DLUTMuted] (0,0) -- (5.7,0);
        \draw[->,DLUTMuted] (0,0) -- (0,106);
        \foreach \y in {0,20,40,60,80,100} {
          \draw[DLUTLine] (0,\y) -- (5.55,\y);
          \node[left,text=DLUTMuted,font=\scriptsize] at (0,\y) {\y};
        }
        \fill[DLUTLine] (0.65,0) rectangle (1.45,76);
        \fill[DLUTBlue!55] (2.35,0) rectangle (3.15,84);
        \fill[DLUTNavy] (4.05,0) rectangle (4.85,92);
        \node[below,text=DLUTMuted,font=\scriptsize] at (1.05,0) {Baseline};
        \node[below,text=DLUTMuted,font=\scriptsize] at (2.75,0) {Method A};
        \node[below,text=DLUTMuted,font=\scriptsize] at (4.45,0) {本文方法};
        \node[above,text=DLUTDeep,font=\scriptsize\bfseries] at (1.05,76) {76.0};
        \node[above,text=DLUTDeep,font=\scriptsize\bfseries] at (2.75,84) {84.0};
        \node[above,text=DLUTDeep,font=\scriptsize\bfseries] at (4.45,92) {92.0};
      \end{tikzpicture}
    \end{column}
    \begin{column}{.42\textwidth}
      \vspace{.1cm}
      \DLUTMetric{准确率提升}{+16.0\%}
      \vspace{.16cm}

      \DLUTMetric{平均响应时间}{38 ms}
      \vspace{.28cm}

      \begin{itemize}
        \item 本文方法在主要指标上表现最佳；
        \item 增加的系统开销处于可接受范围。
      \end{itemize}
    \end{column}
  \end{columns}
\end{frame}

\begin{frame}{消融实验}
  \begin{table}
    \centering
    \small
    \begin{tabular}{lccc}
      \toprule
      方法配置 & Accuracy (\%) & F1-score (\%) & 延迟 (ms) \\
      \midrule
      基础模型 & 76.0 & 74.8 & 24 \\
      + 特征增强 & 84.0 & 82.9 & 31 \\
      + 优化策略（本文方法） & \textbf{92.0} & \textbf{91.4} & 38 \\
      \bottomrule
    \end{tabular}
  \end{table}
  \vspace{.3cm}
  \begin{exampleblock}{结论}
    各模块均带来稳定收益，优化策略对最终效果提升最为明显。
  \end{exampleblock}
\end{frame}

\section{总结与展望}

\begin{frame}{总结与展望}
  \begin{columns}[T,onlytextwidth]
    \begin{column}{.47\textwidth}
      \begin{block}{工作总结}
        \begin{itemize}
          \item 完成需求分析与总体方案设计；
          \item 实现核心方法和系统原型；
          \item 通过多组实验验证方案有效性。
        \end{itemize}
      \end{block}
    \end{column}
    \begin{column}{.47\textwidth}
      \begin{exampleblock}{后续展望}
        \begin{itemize}
          \item 扩充更复杂场景的数据覆盖；
          \item 进一步降低在线推理开销；
          \item 加强可解释性与异常分析能力。
        \end{itemize}
      \end{exampleblock}
    \end{column}
  \end{columns}
\end{frame}

\begin{frame}[plain]
  \begin{tikzpicture}[remember picture,overlay]
    \fill[DLUTDeep] (current page.south west) rectangle (current page.north east);
    \fill[DLUTBlue] (current page.south west) rectangle
      ($(current page.south east)+(0,.16cm)$);
    \node[anchor=north east,opacity=.11] at
      ($(current page.north east)+(-.5cm,-.5cm)$) {\DLUTLogo[5.25cm]};
    \node[text=white,font=\bfseries\fontsize{28}{34}\selectfont] at
      ($(current page.center)+(0,.45cm)$) {感谢各位老师};
    \node[text=DLUTSky,font=\large] at
      ($(current page.center)+(0,-.42cm)$) {敬请批评指正};
    \node[text=DLUTSky,font=\small] at
      ($(current page.center)+(0,-1.42cm)$) {Q \& A};
  \end{tikzpicture}
\end{frame}

\end{document}
